\documentclass{TDP003mall}

\usepackage{float}
\usepackage{minted}


\newcommand{\version}{Version 1.0}
\author{Otto Roming, \url{ottro821@student.liu.se}}
\title{Referat}
\date{2026-03-01}
\rhead{Otto Roming}

\begin{document}
\projectpage
\section{Revisionshistorik}
\begin{table}[!h]
\begin{tabularx}{\linewidth}{|l|X|l|}
\hline
Ver. & Revisionsbeskrivning & Datum \\\hline
1.0 & Första version av referatet & 2025-03-01 \\\hline
\end{tabularx}
\end{table}

\section{Referat}
I \cite{backus_can_1978} beskriver Backus ett fenomen som han kallar för Von Neumann programspråk. Dessa programspråk är byggda för att köras på Von Neumann datorer, som karaktäriseras av att bestå av processorer och lagring där endast ett ord av data kan fraktas mellan processorn och lagringen åt gången. Detta ger dessa system som en flaskhals där endast ett ord av data kan skickas mellan processorn och lagringen vilket saktar ner arbete betydligt.

Backus \cite{backus_can_1978} anser att de flesta konventionella språk så som bland annat \textit{Fortran} och \textit{Algol} är Von Neumann språk. Backus \cite{backus_can_1978} kritiserar dessa Von Neumann programspråk för bland annat att vara komplexa, oflexibla, mångordiga i programkoden, och sakna användbara matematiska operatorer. En av de större nackdelarna med Von Neumann programspråk menar Backus \cite{backus_can_1978} är att de gör skillnad på uttryck som gör beräkningar och satser som hanterar kontrollflöde och tilldelningar, Backus \cite{backus_can_1978} menar med detta att uppdelningen lägger till onödig komplexitet i ett system som hade kunnat vara betydligt enklare.

Backus\cite{backus_can_1978} använder ett kodexempel på Von Neumann språk (se figur \ref{fig:vonneumann}) där han påpekar att läsaren av programkoden måste mentalt exekvera koden för att förstå den, koden är ickegenerell på grund av användningen av a, b, och c vilket innebär att den skulle behöva omringas av en procedurdeklaration för att återanvändas och att koden är begränsad till att behöva använda sig av ett ord åt gången repetition för att modifiera värdet till det önskade.

\begin{figure}[H]
    \centering
    \begin{minipage}{0.5\linewidth}
    \begin{minted}[frame=single]{text}
c = 0
for i = 1 step 1 until n do
    c = c + a[i] × b[i]
    \end{minted}
    \end{minipage}
    \caption{Backus kodexempel på Von Neumann språk \cite{backus_can_1978}}
    \label{fig:vonneumann}
\end{figure}

Backus \cite{backus_can_1978} förespråkar funktionella programspråk och formellt funktionella programspråk som inte använder sig av variabler och istället är uppbyggda av objekt, funktioner och definitioner. Enligt Backus \cite{backus_can_1978} så är blir dessa språk mycket lättare för kodaren att skriva då de inte involverar samma mängd komplexitet som Von Neumann språk oftast gör. Detta gör dem enligt Backus \cite{backus_can_1978} det tydliga svaret till framtidens problem.

Funktionella programspråk har länge blivit kritiserade för att det kan bli komplicerat att spara data mellan körningar av programmet \cite{backus_can_1978}. Här menar Backus \cite{backus_can_1978} att detta problem kan lösas istället med hjälp av ett system som kallas AST (\textit{Applicative State Transition}) vilket används för att enkelt spara programmets tillstånd samt regler för att byta mellan olika tillstånd, viktigt att notera är att detta tillstånd ändras endast efter större beräkningar till skillnad från Von Neumann system som ändrar tillstånd mycket ofta. Backus \cite{backus_can_1978} föreställer sig att användning av dessa AST-system kan ge möjlighet till innovationer för nya datorsystem på grund av att dessa system är utformade efter att kunna utnyttja parallellism till fullt.


% Can programming be liberated from the von neumann style? a functional style and its algebra of programs




\bibliographystyle{IEEEtran}
\bibliography{references}

\end{document}